\Askhsh{25257_SOLUTION}{1}
\begin{ylh}{1}
    \item [ΛΥΣΗ]
\end{ylh}
\begin{alist}
\item Α) Επειδή το Μ διαγράφει το ευθύγραμμο τμήμα ΑΝ με $AM \ne 0$, ισχύει:
\[ 0 < AM \le AN \Leftrightarrow 0 < x \le 0,8 + 1 \Leftrightarrow x \in \left(0, \frac{9}{5}\right]. \]
\begin{rlist}
\item Όταν το Μ κινείται εσωτερικά του ΑΡ, είναι $x \in \left(0, \frac{4}{5}\right]$ και $E=(\text{AKΛA})$. Τα όμοια τρίγωνα ΑΜΛ και ΑΡΕ, έχουν λόγο ομοιότητας $\lambda$
\begin{wrapfigure}[14]{r}{5cm} % Adjust number of lines and width as needed
\begin{tikzpicture}[scale=0.7]
    \tkzDefPoint(0,4){A}
    \tkzDefPoint(-1.5,2){K_diag} % K in diagram
    \tkzDefPoint(1.5,2){P_diag} % P in diagram
    \tkzDefPoint(-2.5,0){B_diag} % B in diagram
    \tkzDefPoint(2.5,0){E_diag} % E in diagram

    \tkzDrawPolygon(A,K_diag,B_diag,E_diag,P_diag)
    \tkzDrawSegment(K_diag,P_diag)

    \tkzLabelPoint[above](A){A}
    \tkzLabelPoint[left](K_diag){K}
    \tkzLabelPoint[right](P_diag){P}
    \tkzLabelPoint[below left](B_diag){B}
    \tkzLabelPoint[below right](E_diag){E}

    % Points M and Lambda on AK_diag and AP_diag respectively
    \tkzDefPointBy[homothety=center A ratio 0.6](K_diag) \tkzGetPoint{M_label}
    \tkzDefPointBy[homothety=center A ratio 0.6](P_diag) \tkzGetPoint{L_label}
    \tkzLabelPoint[left](M_label){M}
    \tkzLabelPoint[right](L_label){$\Lambda$}
    \tkzDrawSegment[dotted](M_label,L_label)
\end{tikzpicture}
\end{wrapfigure}
\[ \text{και είναι } \lambda = \frac{\text{AM}}{\text{AP}} = \frac{x}{0,8} = \frac{5x}{4}. \]
Οπότε ο λόγος των εμβαδών τους είναι:
\[ \frac{(\text{AM}\Lambda)}{(\text{APE})} = \lambda^2 = \left(\frac{5x}{4}\right)^2 = \frac{25x^2}{16}. \]
Έτσι έχουμε:
\begin{align*} \frac{E}{\text{PE} \cdot \text{AP}} &= \frac{25x^2}{16} \\ \Leftrightarrow \frac{E}{0,8 \cdot 0,8} &= \frac{25x^2}{16} \\ \Leftrightarrow \frac{E}{0,64} &= \frac{25x^2}{16} \\ \Leftrightarrow E &= \frac{25x^2}{16} \cdot 0,64 \\ \Leftrightarrow E &= \frac{25x^2}{16} \cdot \frac{64}{100} \\ \Leftrightarrow E &= \frac{25x^2}{16} \cdot \frac{16}{25} \\ \Leftrightarrow E &= x^2. \end{align*}
Σχόλιο: Εναλλακτικά για τον υπολογισμό του εμβαδού $E = (\text{AKΛA})$ από την ομοιότητα των τριγώνων ΑΜΛ και ΑΡΕ, με λόγο ομοιότητας $\lambda$ είναι:
\[ \lambda = \frac{\text{AM}}{\text{AP}} = \frac{x}{0,8} \Leftrightarrow \frac{\text{ΚΛ}}{\text{PE}} = \frac{x}{0,8} \Leftrightarrow \text{ΚΛ} = x. \]
Άρα ΚΛ = $2x$ και $E = \frac{\text{ΚΛ} \cdot \text{AM}}{2} = \frac{2x \cdot x}{2} = x^2 \text{ m}^2$.

\item Όταν το Μ κινείται από το P ως το N, είναι $x \in \left[\frac{4}{5}, \frac{9}{5}\right]$ και $E=(\text{ABE}) + (\text{BKAE})$.
\begin{wrapfigure}[7]{r}{5cm} % Adjust number of lines and width as needed
\begin{tikzpicture}[scale=0.7]
    \tkzDefPoint(0,4){A}
    \tkzDefPoint(-1.5,2){K_diag}
    \tkzDefPoint(1.5,2){P_diag}
    \tkzDefPoint(-2.5,0){B_diag}
    \tkzDefPoint(2.5,0){E_diag}

    \tkzDrawPolygon(A,K_diag,B_diag,E_diag,P_diag)
    \tkzDrawSegment(K_diag,P_diag)

    \tkzLabelPoint[above](A){A}
    \tkzLabelPoint[left](K_diag){K}
    \tkzLabelPoint[right](P_diag){P}
    \tkzLabelPoint[below left](B_diag){B}
    \tkzLabelPoint[below right](E_diag){E}

    % M is the foot of the altitude from A to KP.
    \tkzDefMidPoint(K_diag,P_diag) \tkzGetPoint{M_altitude_foot}
    \tkzDrawSegment[dashed](A,M_altitude_foot)
    \tkzLabelPoint[below](M_altitude_foot){M}
    
    % Vertical line from M_altitude_foot to B_diagE_diag
    \tkzDefPointWith[projection](M_altitude_foot,B_diag,E_diag) \tkzGetPoint{H_base}
    \tkzDrawSegment[dashed](M_altitude_foot,H_base)
\end{tikzpicture}
\end{wrapfigure}
Έτσι, είναι:
\begin{align*} E &= \frac{\text{BE} \cdot \text{AP}}{2} + \text{BE} \cdot \text{PM} \\ &= \frac{0,8 \cdot 1,6}{2} + 1,6 \cdot (x-0,8) \\ &= 0,8 \cdot 0,8 + 1,6 \cdot (x-0,8) \\ &= \frac{4}{5} \cdot \frac{4}{5} + \frac{8}{5} \cdot \left(x-\frac{4}{5}\right) \\ &= \frac{16}{25} + \frac{8}{5}x - \frac{32}{25} \\ &= \frac{8}{5}x - \frac{16}{25} \text{ m}^2. \end{align*}
\end{rlist}
Οπότε το εμβαδόν του παραθύρου που καλύπτει η σίτα σε $\text{m}^2$ είναι:
\[ E(x) = \begin{cases} x^2, & \text{αν } x \in \left(0, \frac{4}{5}\right] \\ \frac{8}{5}x - \frac{16}{25}, & \text{αν } x \in \left[\frac{4}{5}, \frac{9}{5}\right] \end{cases} \]

\item β) Είναι $E\left(\frac{4}{5}\right) = \frac{8}{5} \cdot \frac{4}{5} - \frac{16}{25} = \frac{32}{25} - \frac{16}{25} = \frac{16}{25}$.
έτσι έχουμε:
\begin{align*} \lim_{x \to \frac{4}{5}^-} E(x) &= \lim_{x \to \frac{4}{5}^-} x^2 = \left(\frac{4}{5}\right)^2 = \frac{16}{25}. \\ \text{και } \lim_{x \to \frac{4}{5}^+} E(x) &= \lim_{x \to \frac{4}{5}^+} \left(\frac{8}{5}x - \frac{16}{25}\right) = \frac{8}{5} \cdot \frac{4}{5} - \frac{16}{25} = \frac{32}{25} - \frac{16}{25} = \frac{16}{25}. \end{align*}
Επειδή $\lim_{x \to \frac{4}{5}^-} E(x) = \lim_{x \to \frac{4}{5}^+} E(x) = E\left(\frac{4}{5}\right) = \frac{16}{25}$, έπεται ότι $\lim_{x \to \frac{4}{5}} E(x) = \frac{16}{25}$.
Άρα η ζητούμενη συνάρτηση ρυθμός μεταβολής είναι ο $E'\left(\frac{4}{5}\right) \frac{8}{5} \text{ m}^2/\text{m}$.

\item γ) Αν είναι $x'(t_0) = 0,08 \text{ m/s}$, για μια οποιαδήποτε χρονική στιγμή $t_0$ τότε ισχύουν:
\begin{itemize}
\item $x(t_0) = \frac{4}{5} \text{ m}$,
\item $E'\left(\frac{4}{5}\right) \text{ m}^2/\text{m}$ και
\item $E\left(\frac{4}{5}\right) \text{ m}^2$.
\end{itemize}
Οπότε ο ρυθμός μεταβολής του εμβαδού Ε τη χρονική στιγμή $t_0$ είναι ίσος με:
\[ (E \circ x)'(t_0) = E'(x(t_0)) \cdot x'(t_0) = E'\left(\frac{4}{5}\right) \cdot \frac{2}{25} = \frac{16}{25} \text{ m}^2/\text{s}. \]
\end{alist}